\documentclass[12pt,a4paper]{article}
\usepackage[utf8]{inputenc}
\usepackage[portuguese]{babel}
\usepackage{amsmath,amssymb,amsfonts}
\usepackage{geometry}
\usepackage{xcolor}
\usepackage{tcolorbox}
\usepackage{tikz}
\usepackage{booktabs}
\usepackage{hyperref}

\geometry{
  top=2.5cm,
  bottom=2.5cm,
  left=2.5cm,
  right=2.5cm
}

\definecolor{primaryblue}{RGB}{24, 76, 120}
\definecolor{accentorange}{RGB}{204, 85, 0}
\definecolor{bgcode}{RGB}{245, 247, 250}

\tcbset{
  colback=bgcode,
  colframe=primaryblue,
  fonttitle=\bfseries,
  arc=3mm
}

\title{\textbf{\huge Compreendendo o \textit{Drift} em Transformações Afins 2D}\\[0.5em]
\large A Origem Geométrica, o Sanduíche de Pivô e a Invariância Espacial em Motores Gráficos}
\author{\textbf{anicrop Engine} — Documentação Teórica}
\date{\today}

\begin{document}

\maketitle
\tableofcontents
\newpage

\section{Introdução: O que é o \textit{Drift}?}

No contexto de motores de composição gráfica 2D (como o \texttt{anicrop}, Photoshop, Figma ou After Effects), o termo \textbf{\textit{Drift}} refere-se a:

\begin{tcolorbox}[title=Definição Formal de \textit{Drift}]
O deslocamento angular e translacional espúrio que surge na caixa delimitadora (\textit{Bounding Box}) de uma camada quando uma transformação afim (rotação, escala ou cisalhamento) é aplicada em torno de um ponto de pivô, fazendo com que os cantos transformados não comecem na origem $(0, 0)$ do espaço local.
\end{tcolorbox}

O \textit{drift} \textbf{não é um erro de arredondamento de ponto flutuante} nem um bug de código. Ele é uma consequência geométrica e trigonométrica fundamental da projeção de retângulos rotacionados em coordenadas cartesianas discretas.

---

\section{A Álgebra das Matrizes Afins e o Sanduíche de Pivô}

\subsection{Representação por Coordenadas Homogêneas $3 \times 3$}

Em computação gráfica 2D, qualquer transformação afim sobre um ponto $\mathbf{p} = \begin{bmatrix} x \\ y \end{bmatrix}$ é expressa na forma homogênea:
\begin{equation}
\begin{bmatrix} x' \\ y' \\ 1 \end{bmatrix} = \mathbf{M} \begin{bmatrix} x \\ y \\ 1 \end{bmatrix} = \begin{bmatrix} m_{00} & m_{01} & t_x \\ m_{10} & m_{11} & t_y \\ 0 & 0 & 1 \end{bmatrix} \begin{bmatrix} x \\ y \\ 1 \end{bmatrix}
\end{equation}

Onde a submatriz $2 \times 2$ superior esquerda controla rotação, escala e cisalhamento (\textit{distorção}), e o vetor $\begin{bmatrix} t_x \\ t_y \end{bmatrix}$ controla a translação.

\subsection{O Sanduíche de Pivô (\textit{Pivot Sandwich})}

A rotação pura padrão $\mathbf{R}(\theta)$ gira em torno da origem $(0, 0)$:
\begin{equation}
\mathbf{R}(\theta) = \begin{bmatrix} \cos\theta & -\sin\theta & 0 \\ \sin\theta & \cos\theta & 0 \\ 0 & 0 & 1 \end{bmatrix}
\end{equation}

Para girar uma imagem de tamanho $W \times H$ em torno de um pivô arbitrário $\mathbf{P} = (p_x, p_y)$ (por exemplo, seu centro $(\frac{W}{2}, \frac{H}{2})$), o motor gráfico deve:
\begin{enumerate}
    \item Transladar o pivô $\mathbf{P}$ para a origem: $\mathbf{T}(-\mathbf{P}) = \begin{bmatrix} 1 & 0 & -p_x \\ 0 & 1 & -p_y \\ 0 & 0 & 1 \end{bmatrix}$
    \item Aplicar a rotação pura: $\mathbf{R}(\theta)$
    \item Transladar a origem de volta para o pivô: $\mathbf{T}(\mathbf{P}) = \begin{bmatrix} 1 & 0 & p_x \\ 0 & 1 & p_y \\ 0 & 0 & 1 \end{bmatrix}$
\end{enumerate}

A matriz resultante da rotação com pivô é o produto:
\begin{equation}
\mathbf{M}_{\text{pivô}} = \mathbf{T}(\mathbf{P}) \cdot \mathbf{R}(\theta) \cdot \mathbf{T}(-\mathbf{P})
\end{equation}

Expandindo algebricamente:
\begin{equation}
\mathbf{M}_{\text{pivô}} = \begin{bmatrix} 
\cos\theta & -\sin\theta & p_x(1 - \cos\theta) + p_y\sin\theta \\
\sin\theta & \cos\theta & p_y(1 - \cos\theta) - p_x\sin\theta \\
0 & 0 & 1
\end{bmatrix}
\end{equation}

\begin{tcolorbox}[title=A Origem Algébrica do Drift]
Observe o vetor de translação introduzido automaticamente pelo sanduíche de pivô:
\begin{align*}
t_x(\mathbf{P}, \theta) &= p_x(1 - \cos\theta) + p_y\sin\theta \\
t_y(\mathbf{P}, \theta) &= p_y(1 - \cos\theta) - p_x\sin\theta
\end{align*}
Esse vetor \textbf{muda drasticamente} conforme o tamanho da imagem ou a posição do pivô varia!
\end{tcolorbox}

---

\section{A Geometria dos 4 Cantos e a Bounding Box (AABB)}

Considere um retângulo alinhado aos eixos com vértices:
\begin{equation}
C_0 = (0, 0), \quad C_1 = (W, 0), \quad C_2 = (W, H), \quad C_3 = (0, H)
\end{equation}

Quando multiplicamos esses 4 vértices pela matriz $\mathbf{M}_{\text{pivô}}$, obtemos 4 novos pontos no espaço 2D: $C'_0, C'_1, C'_2, C'_3$.

A \textit{Axis-Aligned Bounding Box} (AABB) resultante no Canvas é definida pelos extremos:
\begin{align}
X_{\min} &= \min(C'_{0x}, C'_{1x}, C'_{2x}, C'_{3x}) \\
Y_{\min} &= \min(C'_{0y}, C'_{1y}, C'_{2y}, C'_{3y}) \\
X_{\max} &= \max(C'_{0x}, C'_{1x}, C'_{2x}, C'_{3x}) \\
Y_{\max} &= \max(C'_{0y}, C'_{1y}, C'_{2y}, C'_{3y})
\end{align}

\subsection{O Vetor de Drift Intrínseco: $(\Delta X_{\text{drift}}, \Delta Y_{\text{drift}})$}

O vetor de \textit{drift} é exatamente a coordenada do canto superior esquerdo da Bounding Box no espaço transformado local:
\begin{equation}
\mathbf{Drift}(\mathbf{M}, W, H) = (X_{\min}, Y_{\min})
\end{equation}

\textbf{Exemplo Numérico Prático ($W=400, H=400, \theta=45^\circ$):}
\begin{itemize}
    \item Pivô central: $\mathbf{P} = (200, 200)$
    \item $\cos(45^\circ) = \sin(45^\circ) = \frac{\sqrt{2}}{2} \approx 0.7071$
    \item Os cantos transformados são:
    \begin{align*}
    C'_0 &= (200 - 200\sqrt{2}, 200) \approx (-82.84, 200) \\
    C'_1 &= (200, 200 - 200\sqrt{2}) \approx (200, -82.84) \\
    C'_2 &= (200 + 200\sqrt{2}, 200) \approx (482.84, 200) \\
    C'_3 &= (200, 200 + 200\sqrt{2}) \approx (200, 482.84)
    \end{align*}
    \item $X_{\min} = -82.84 \approx -83\text{px}$
    \item $Y_{\min} = -82.84 \approx -83\text{px}$
    \item Largura da Bounding Box: $482.84 - (-82.84) = 400\sqrt{2} \approx 565.68 \approx 566\text{px}$.
\end{itemize}

O \textit{drift} dessa rotação é $\mathbf{Drift} = (-83, -83)$. Se você posicionar a moldura em $(X, Y)$, a Bounding Box no Canvas começará em $(X - 83, Y - 83)$!

---

\section{O Caso Crítico: Crop, Rotação e Desdobramento (\textit{Uncrop / Fit Content})}

Aqui está a anatomia exata do problema que investigamos no cenário Mariachi.

\subsection{Cenário 1: Camada Cropada e Rotacionada}
\begin{enumerate}
    \item Imagem base: $736 \times 1104$.
    \item Recorte (ROI): $(100, 50, 400, 400)$.
    \item Rotação: $\theta = 45^\circ$. O pivô é o centro do recorte: $\mathbf{P}_{\text{crop}} = (200, 200)$.
    \item Alinhamento ao Canvas: A moldura vai para $(168, 352)$.
    \item O \textit{offset} da base em relação ao recorte é:
    $$\mathbf{Offset} = \mathbf{Base}_{\text{pos}} - \mathbf{Layout}_{\text{pos}} = (68, 302) - (168, 352) = (-100, -50)$$
    \item A matriz física que renderiza os pixels da imagem no Canvas é:
    \begin{equation}
    \mathbf{M}_{\text{content}} = \mathbf{T}(168, 352) \cdot \mathbf{M}_{\text{pivô}}(400, 400, 45^\circ) \cdot \mathbf{T}(-100, -50)
    \end{equation}
    Calculando numericamente a translação dos pixels no Canvas:
    $$\mathbf{M}_{\text{content}}[0, 2] = 332.64\text{px}, \quad \mathbf{M}_{\text{content}}[1, 2] = 163.09\text{px}$$
\end{enumerate}

\subsection{Cenário 2: O Que Acontecia no \textit{Fit Content} Sem Compensação}

Ao chamar \texttt{fit\_content}, o recorte é desfeito. O tamanho de referência passa de $400 \times 400$ para $736 \times 1104$.

\begin{enumerate}
    \item O novo pivô na imagem base passa a ser: $\mathbf{P}_{\text{base}} = (368, 552)$.
    \item O sanduíche de pivô para a base calcula um novo vetor de translação intrínseca:
    $$t_x(\mathbf{P}_{\text{base}}, 45^\circ) \approx 498.14, \quad t_y(\mathbf{P}_{\text{base}}, 45^\circ) \approx -139.19$$
    \item Se a moldura for colocada em $(132, 246)$, mas a base continuar com $\mathbf{Offset} = (-64, +56)$, a matriz de renderização dos pixels é corrompida:
    $$\mathbf{M}_{\text{novo}}[0, 2] = 247.15\text{px} \quad (\Delta X = 332.64 - 247.15 = \mathbf{85.5\text{px}})$$
    $$\mathbf{M}_{\text{novo}}[1, 2] = 157.50\text{px} \quad (\Delta Y = 163.09 - 157.50 = \mathbf{5.6\text{px}})$$
\end{enumerate}

\textbf{Resultado Visual:} A imagem no GIMP salta $85.5\text{px}$ para o lado, quebrando o alinhamento visual!

---

\section{A Solução Algébrica Definitiva: Invariância via Matriz Inversa}

Para que nenhum pixel se mova no Canvas durante uma mudança de enquadramento, impomos a \textbf{Condição de Invariância Global}:
\begin{equation}
\mathbf{M}_{\text{content\_depois}} \equiv \mathbf{M}_{\text{content\_antes}}
\end{equation}

Como a nova camada deve operar sob a matriz de distorção da nova referência $\mathbf{M}_{\text{transform}}$:
\begin{equation}
\mathbf{T}(x, y) \cdot \mathbf{M}_{\text{transform}} = \mathbf{M}_{\text{content\_alvo}}
\end{equation}

Multiplicando à direita pela matriz inversa $\mathbf{M}_{\text{transform}}^{-1}$:
\begin{equation}
\mathbf{T}(x, y) = \mathbf{M}_{\text{content\_alvo}} \cdot \mathbf{M}_{\text{transform}}^{-1}
\end{equation}

\begin{tcolorbox}[title=Propriedade Fundamental da Inversão Afim]
Como ambas as matrizes compartilham o mesmo bloco de rotação $\mathbf{R}(45^\circ)$, o produto $\mathbf{R}(45^\circ) \cdot \mathbf{R}(-45^\circ) = \mathbf{I}$ anula a rotação completamente!
O produto resultante é uma \textbf{Matriz de Translação Pura}:
\begin{equation}
\begin{bmatrix} 1 & 0 & x \\ 0 & 1 & y \\ 0 & 0 & 1 \end{bmatrix}
\end{equation}
Onde $x$ e $y$ são exatamente as coordenadas onde a nova moldura deve ser ancorada para zerar o \textit{drift} e o \textit{offset} residual.
\end{tcolorbox}

---

\section{Resumo dos Princípios Arquiteturais Consolidados}

\begin{table}[h!]
\centering
\begin{tabular}{@{}lll@{}}
\toprule
\textbf{Conceito} & \textbf{Responsabilidade} & \textbf{Regra de Ouro} \\ \midrule
\texttt{Transform} & Distorção Afim Intrínseca & Sanduíche de pivô local imutável \\
\texttt{Layout} & Enquadramento e Posição & Compensa o drift angular via $\mathbf{T} = \mathbf{M} \cdot \mathbf{M}_{\text{trans}}^{-1}$ \\
\texttt{BaseLayer.\_region} & Geometria Canônica da Imagem & Sincronizada ao zerar o crop (\textit{offset} $= 0$) \\
\texttt{Content Matrix} & Projeção dos Pixels no Canvas & Imutável durante operações puras de layout \\ \bottomrule
\end{tabular}
\caption{Separação de responsabilidades no motor geométrico do \texttt{anicrop}.}
\end{table}

\end{document}
